% Owner: Kashaf Ali / Member A. Complete rewrite authorized 5 October 2026.
% Local font encoding restores the template's Times regular/bold faces under XeTeX.
\begingroup
\renewcommand{\encodingdefault}{T1}
\fontencoding{T1}\selectfont
\chapter{Introduction}

Clinical decision support brings together patient information, documentation and analytical assistance. The following sections establish the background, purpose, intended readership and terminology needed to understand the work, followed by an outline of the remaining chapters.

\section{Purpose of this Document}

Zeest investigates how longitudinal memory, agent coordination, evidence retrieval and critic review can support clinical decision making. The central comparison examines a single language model, a memory-augmented variant and the full system, with controlled contrasts for critic and multimodal contributions. The requirements and design establish the physician workflow, component responsibilities and information relationships needed for these comparisons. Existing validated diagnostic models form the supporting groundwork; the research focuses on their integration with memory and review. Clinical coverage depends on suitable data and models, and evaluation uses de-identified or synthetic cases. The intended outcome is an assessable decision-support system whose recommendations remain subject to physician judgement. The design also distinguishes authorised clinical observations from generated drafts, identifies unresolved resource choices and establishes source-aware review. Implementation and comparative evaluation remain future work, so the described workflow sets requirements without presenting unmeasured outcomes or claiming clinical deployment. The assessment requires defined comparison conditions.

\section{Intended Audience}

The primary audience includes the project supervisor, academic evaluators and the student development team. Evaluators need a clear account of the clinical problem, research contribution and technical decisions, while developers need consistent requirements and component responsibilities. Physicians are also relevant readers because Zeest centres on their interpretation of patient information and review of generated recommendations. The descriptions of findings, evidence and approval controls should therefore remain understandable without requiring knowledge of agent orchestration. Researchers interested in clinical language models, memory and retrieval can examine how these capabilities are integrated. The writing connects technical responsibilities to user interactions so that each audience can follow the integrated system's intended behaviour. Interface designers can use the screen mappings to connect user tasks with requirements, while database and integration developers can check ownership and revision relationships. Clear distinctions between design, source evidence and measured outcomes help these readers assess the same clinical workflow consistently.

\section{Definitions, Acronyms, and Abbreviations}

Key terms and abbreviations are defined below.

\noindent\textbf{AI:} Artificial Intelligence.\\
\noindent\textbf{LLM:} Large Language Model.\\
\noindent\textbf{CDSS:} Clinical Decision Support System.\\
\noindent\textbf{RAG:} Retrieval-Augmented Generation.\\
\noindent\textbf{SDG:} Sustainable Development Goal.\\
\noindent\textbf{Encounter:} An identifiable consultation or visit associated with a patient.\\
\noindent\textbf{Longitudinal memory:} Patient-specific information retained and retrieved across encounters.\\
\noindent\textbf{Hybrid memory:} Relational records combined with semantic vector representations.\\
\noindent\textbf{Multimodal input:} Clinical information represented through notes, images, laboratory information and medication context.\\
\noindent\textbf{Agent:} A specialised processing component with a defined responsibility.\\
\noindent\textbf{Clinical draft:} A generated diagnosis and treatment plan awaiting physician review.\\
\noindent\textbf{Approved decision:} Clinical content explicitly accepted by the physician for recording.\\
\noindent\textbf{Audit history:} Recorded physician actions associated with review and decision making.\\
\noindent\textbf{Ablation study:} A comparison of system variants used to examine component contributions.


% Finish this section's floats before the next heading.

\section{Conclusion}

Chapter~1 introduces the information problem and the purpose of Zeest. Chapter~2 defines the vision, objectives, scope and stakeholder needs. Chapter~3 examines the related research and applications that inform the technical approach. Chapter~4 specifies features, functional and non-functional requirements, use cases, resources, interface designs, data relationships and risks. Chapter~5 develops the methodology and comparative evaluation approach. Chapter~6 presents the high-level and low-level design of the integrated workflow. References identify the supporting sources, and the Appendix contains supplementary material. Together, these chapters connect the clinical motivation to a system design and an evaluation framework for investigating memory, reasoning, critic review and multimodal contributions. The resulting structure preserves the distinction between recorded observations, generated recommendations and accepted decisions. Source-aware records and review actions establish the connection between clinical assistance and accountability. Subsequent implementation must preserve these relationships before comparative evaluation can establish the value of each component.

\endgroup
